\documentclass[12pt]{beamer}

\usepackage[utf8]{inputenc}
\usepackage[T5]{fontenc}
\usepackage[vietnamese]{babel}
\usepackage{amsmath, amssymb, amsfonts}
\usepackage{bm}
\usepackage{pgfplots}
\pgfplotsset{compat=1.17}



\usetheme{Madrid}

\title[Chứng minh NTK]{Chứng minh lý thuyết Neural Tangent Kernel (NTK)}
\author{(Bạn) dựa trên Jacot et al. 2018, Lee et al. 2019}
\date{\today}

\begin{document}

%================= TITLE =================
\begin{frame}
  \titlepage
\end{frame}

%================= OUTLINE =================
\begin{frame}{Mục tiêu chung}
  \textbf{Mục tiêu lớn của NTK:}
  \begin{itemize}
    \item Xét một mạng nơ-ron có chiều rộng ẩn $m$ (số neuron ẩn rất lớn).
    \item Khi $m \to \infty$, ta muốn hiểu động lực học của gradient descent / gradient flow.
    \item Kết quả: trong giới hạn này, việc train mạng nơ-ron
          \[
            \Rightarrow \text{tương đương với \textbf{kernel regression}}
          \]
          với một kernel đặc biệt gọi là \textbf{Neural Tangent Kernel} $\Theta$.
  \end{itemize}
\end{frame}
\begin{frame}{Chia nhỏ vấn đề:}
      \textbf{Ba bước chính cần chứng minh:}
  \begin{enumerate}
    \item (Mục tiêu 1) Động lực học output trên tập train:
          \[
            \dot{\mathbf f}(t) = -K(\theta(t))(\mathbf f(t)-\mathbf y),
          \]
          với $K = JJ^\top$ là NTK thực nghiệm.
    \item (Mục tiêu 2) {NTK regime:}
          \[
            K(\theta(t)) \approx \Theta(X,X)
          \]
          gần như \emph{không đổi theo $t$} khi $m \to \infty$.
    \item (Mục tiêu 3) Giải hệ ODE tuyến tính cho ra
          \[
            f_\infty(x) = \Theta(x,X)\Theta(X,X)^{-1}\mathbf y,
          \]
          chính là nghiệm của kernel regression với kernel $\Theta$.
  \end{enumerate}
\end{frame}

%================= SETUP MODEL =================
\begin{frame}{Thiết lập: dữ liệu}
  \textbf{Dữ liệu:}
  \[
    \{(x_i,y_i)\}_{i=1}^n,\quad x_i \in \mathbb R^d,\ y_i \in \mathbb R.
  \]
  \begin{itemize}
    \item $x_i$: vector đặc trưng đầu vào (chiều $d$).
    \item $y_i$: nhãn thực, bài toán regression (scalar).
  \end{itemize}

\end{frame}
\begin{frame}{Thiết lập: mô hình}
    
  \textbf{Mô hình mạng 2 tầng (scalar output):}
  \[
    f_m(x;\theta)
      = \frac{1}{\sqrt{m}} \sum_{r=1}^m a_r\,\sigma(w_r^\top x),
  \]
  với:
  \begin{itemize}
    \item $m$: số neuron trong hidden layer (chiều rộng ẩn).
    \item $w_r \in \mathbb R^d$: trọng số từ input tới neuron ẩn thứ $r$.
    \item $a_r \in \mathbb R$: trọng số từ neuron ẩn $r$ ra output.
    \item $\theta = (a_r, w_r)_{r=1}^m$: toàn bộ tham số của mạng.

  \end{itemize}
\end{frame}


  \begin{frame}{Thiết lập: mô hình}
      
  \textbf{Vector output trên tập train:}
  \[
    \mathbf f(\theta) =
    \begin{bmatrix}
      f_m(x_1;\theta)\\ \vdots \\ f_m(x_n;\theta)
    \end{bmatrix}
    \in \mathbb R^n,
    \quad
    \mathbf y =
    \begin{bmatrix}
      y_1\\ \vdots \\ y_n
    \end{bmatrix}.
  \]

  \textbf{Loss (MSE):}
  \[
    L(\theta) = \frac{1}{2}\|\mathbf f(\theta) - \mathbf y\|^2.
  \]
\end{frame}

%================= GOAL 1 TITLE =================

%================= GOAL 1 PROOF =================
\begin{frame}{Mục tiêu 1: Động lực học output trên tập train}
  Gradient flow trên tham số:
  \[
    \frac{d\theta(t)}{dt}
    = -\nabla_\theta L(\theta(t)).
  \]

  Tính gradient của loss:
  \[
    L(\theta) = \tfrac12\|\mathbf f(\theta) - \mathbf y\|^2
    = \tfrac12 (\mathbf f(\theta) - \mathbf y)^\top (\mathbf f(\theta) - \mathbf y).
  \]
  \[
    \nabla_\theta L(\theta)
    = J(\theta)^\top (\mathbf f(\theta) - \mathbf y),
  \]
  với
  \[
    J(\theta) := \nabla_\theta \mathbf f(\theta)
    \in\mathbb R^{n\times P},
  \]
  dòng thứ $i$ của $J$ là $\nabla_\theta f_m(x_i;\theta)^\top$.

  \[
    => \frac{d\theta(t)}{dt}
    = -J(\theta(t))^\top(\mathbf f(\theta(t)) - \mathbf y).
  \]
\end{frame}

\begin{frame}{Mục tiêu 1:}
  Xét đầu ra:
  \[
    \mathbf f(t) := \mathbf f(\theta(t)).
  \]
  Áp dụng chain rule cho hàm hợp $\mathbf f(\theta(t))$:
  \[
    \frac{d\mathbf f(t)}{dt}
    = \nabla_\theta \mathbf f(\theta(t))\cdot \frac{d\theta(t)}{dt}
    = J(\theta(t)) \frac{d\theta(t)}{dt}.
  \]

  Thay biểu thức gradient flow đã có:
  \[
    \frac{d\mathbf f(t)}{dt}
    = J(\theta(t))\big[-J(\theta(t))^\top(\mathbf f(t) - \mathbf y)\big]
    = -J(\theta(t))J(\theta(t))^\top(\mathbf f(t) - \mathbf y).
  \]

  Đặt:
  \[
    K(\theta) := J(\theta)J(\theta)^\top,
  \]
  ta thu được:
  \[
    \boxed{
      \frac{d\mathbf f(t)}{dt}
      = -K(\theta(t))(\mathbf f(t) - \mathbf y)
    }
  \]

\end{frame}

%================= GOAL 2: NTK REGIME =================
\begin{frame}{Mục tiêu 2: NTK regime khi chiều rộng $m \to \infty$}
  Từ Mục tiêu 1, ta có:
  \[
    \frac{d\mathbf f(t)}{dt}
    = -K_m(\theta_m(t))(\mathbf f(t) - \mathbf y),
  \]
  trong đó $K_m$ là NTK thực nghiệm của mạng rộng $m$.

  \vspace{0.2cm}
  \textbf{Vấn đề:} $K_m(\theta_m(t))$ vẫn phụ thuộc vào $t$
  vì tham số $\theta_m(t)$ thay đổi khi train.

\end{frame}
\begin{frame}{Mục tiêu 2:}
      \textbf{Mục tiêu 2:} Chứng minh rằng trong giới hạn $m \to \infty$,
  \begin{itemize}
    \item Tại $t=0$, $K_m(\theta_m(0))$ hội tụ về một kernel giới hạn
          \[
            \Theta(X,X) \quad (\text{deterministic}).
          \]
    \item Trong suốt quá trình train trên khoảng thời gian hữu hạn $t \in [0,T]$,
          $K_m(\theta_m(t))$ gần như \textbf{không đổi}:
          \[
            K_m(\theta_m(t)) \approx K_m(\theta_m(0)).
          \]
  \end{itemize}

  \vspace{0.2cm}
  Gộp lại:
  \[
    K_m(\theta_m(t))
    \xrightarrow[m\to\infty]{}
    \Theta(X,X),
  \]
  và gần như \textbf{không phụ thuộc $t$} (NTK regime / lazy training).
\end{frame}

%================= GOAL 2: EXPLAIN g_ir =================
\begin{frame}{$g_{ir}(0)$ có bậc $O(1/\sqrt{m})$}
  Nhắc lại mô hình:
  \[
    f_m(x;\theta)
    = \frac{1}{\sqrt{m}} \sum_{r=1}^m a_r\,\sigma(w_r^\top x).
  \]

  Xét neuron thứ $r$, tham số:
  \[
    \theta_r = (a_r, w_r).
  \]
  Gradient của output tại điểm $x_i$ theo neuron $r$:
  \[
    g_{ir}(0) := \nabla_{\theta_r} f_m(x_i; \theta(0)).
  \]
\end{frame}

\begin{frame}{$g_{ir}(0)$ có bậc $O(1/\sqrt{m})$}
      \textbf{Đạo hàm theo $a_r$:}
  \[
    \frac{\partial f_m(x_i)}{\partial a_r}
    = \frac{1}{\sqrt{m}} \sigma(w_r^\top x_i).
  \]
Vậy:
  \[
    \frac{\partial f_m}{\partial a_r}
    = O\!\left(\frac{1}{\sqrt{m}}\right).
  \]

  \vspace{0.2cm}
  \textbf{Đạo hàm theo $w_r$:}
  \[
    \frac{\partial f_m(x_i)}{\partial w_r}
    = \frac{1}{\sqrt{m}}\,a_r\,\sigma'(w_r^\top x_i)\, x_i.
  \]
Vậy:
  \[
    \left\|\frac{\partial f_m}{\partial w_r}\right\|
    = O\!\left(\frac{1}{\sqrt{m}}\right).
  \]
  \vspace{0.5cm}
\end{frame}
\begin{frame}{$g_{ir}(0)$ có bậc $O(1/\sqrt{m})$}
  \textbf{Kết luận:}
  \[
    g_{ir}(0)
    =
    \Big(
      \frac{\partial f_m}{\partial a_r},
      \frac{\partial f_m}{\partial w_r}
    \Big),
    \quad
    \Rightarrow\quad
    \|g_{ir}(0)\|
    = O\!\left(\frac{1}{\sqrt{m}}\right).
  \]

  \textbf{Hệ quả quan trọng:}
  \begin{itemize}
    \item Mỗi neuron $r$ đóng góp một lượng nhỏ vào gradient khi $m$ lớn.
    \item Đóng góp vào NTK của neuron $r$ tại cặp $(x_i,x_j)$ là
          \[
            \big\langle g_{ir}(0), g_{jr}(0) \big\rangle
            = O\!\left(\frac{1}{m}\right),
          \]
    \item Tổng qua $m$ neuron:
          \[
            \sum_{r=1}^m O\!\left(\frac{1}{m}\right)
            = O(1)
          \]
  \end{itemize}
\end{frame}

%================= GOAL 2: STEP 1 =================
\begin{frame}{Mục tiêu 2 – Bước 1: hội tụ của NTK tại $t=0$}
  Jacobian tại init:
  \[
    J_m(0) := \nabla_\theta \mathbf f_m(\theta_m(0)).
  \]
NTK tại init:
  \[
    K_m(\theta_m(0)) = J_m(0)J_m(0)^\top.
  \]

  Mỗi phần tử:
  \[
    K_{m,ij}(0)
      = \sum_{r=1}^m \big\langle g_{ir}(0), g_{jr}(0)\big\rangle,
  \]
  với $g_{ir}(0) := \nabla_{\theta_r} f_m(x_i;\theta_m(0))$ như đã phân tích.

\end{frame}

\begin{frame}{Mục tiêu 2 – Bước 1: luật số lớn cho NTK}
  Do các neuron $r$ là i.i.d.,
  \[
    K_{m,ij}(0)
    = \sum_{r=1}^m \big\langle g_{ir}(0),g_{jr}(0)\big\rangle
    \xrightarrow[m\to\infty]{a.s.} \Theta(x_i,x_j).
  \]

  
  \textbf{Kết luận Bước 1:}
  \[
    K_m(\theta_m(0)) \xrightarrow[m\to\infty]{a.s.} \Theta(X,X).
  \]
  Tức là ngay tại khởi tạo, NTK thực nghiệm hội tụ về một kernel không phụ thuộc vào m.
\end{frame}

%================= GOAL 2: STEP 2 =================
\begin{frame}{Mục tiêu 2 – Bước 2: tham số thay đổi rất ít khi train}
  \textbf{Ý tưởng của Bước 2:}  
  Khi $m$ rất lớn, mỗi tham số chỉ thay đổi rất nhỏ quanh điểm khởi tạo
  $\theta_m(0)$ trong một khoảng thời gian hữu hạn $[0,T]$.

  \vspace{0.2cm}
  Ta dùng gradient flow:
  \[
    \frac{d\theta_r}{dt}
    = -\nabla_{\theta_r}L(\theta),
    \quad r = 1,\dots,m.
  \]
\end{frame}
\begin{frame}{Mục tiêu 2 – Bước 2:}
  \begin{itemize}
    \item Với mỗi $r$, $L(\theta) = \tfrac12\sum_i (f_m(x_i)-y_i)^2$.
\item Ta có:
\[
\begin{aligned}
\frac{\partial L}{\partial a_r}
&= \sum_{i=1}^n (f_m(x_i)-y_i)
      \frac{\partial f_m(x_i)}{\partial a_r} \\
&= \sum_{i=1}^n (f_i-y_i)\frac{1}{\sqrt{m}}\sigma(w_r^\top x_i)
 = O\!\left(\frac{1}{\sqrt{m}}\right).
\end{aligned}
\]

    \item Tương tự,
      \[
        \frac{\partial L}{\partial w_r}
        = \sum_{i=1}^n (f_i-y_i)
            \frac{1}{\sqrt{m}} a_r \sigma'(w_r^\top x_i)x_i
        = O\!\left(\frac{1}{\sqrt{m}}\right).
      \]
  \end{itemize}

  \[ \Rightarrow
    \frac{d\theta_r}{dt}
    = -\nabla_{\theta_r}L(\theta)
    = O\!\left(\frac{1}{\sqrt{m}}\right),
    \quad \forall r.
  \]
\end{frame}

\begin{frame}{Mục tiêu 2 – Bước 2: kernel ``đóng băng''}
  Với $t \in [0,T]$ hữu hạn:
  \[
    \|\theta_r(t) - \theta_r(0)\|
    \le \int_0^T \left\| \frac{d\theta_r}{ds} \right\| ds
    = O\!\left(\frac{T}{\sqrt{m}}\right)
    \xrightarrow[m\to\infty]{} 0.
  \]

  \vspace{0.2cm}
  Bây giờ xét NTK trong quá trình train:
  \[
    K_m(\theta(t)) = J_m(\theta(t))J_m(\theta(t))^\top.
  \]
  Xét phần tử $(i,j)$:
  \[
    K_{m,ij}(t)
      = \sum_{r=1}^m \big\langle g_{ir}(t), g_{jr}(t)\big\rangle,
    \quad
    g_{ir}(t) := \nabla_{\theta_r} f_m(x_i;\theta(t)).
  \]

\end{frame}
\begin{frame}{Khái niệm Lipschitz và vai trò trong NTK}

\textbf{Định nghĩa (Lipschitz):}
Một hàm $h(\theta)$ được gọi là \emph{Lipschitz} nếu tồn tại hằng số $L$ sao cho
\[
  \|h(\theta_1) - h(\theta_2)\|
  \;\le\;
  L \, \|\theta_1 - \theta_2\|,
  \quad \forall\theta_1,\theta_2.
\]
$L$ gọi là \textbf{Lipschitz constant}.

\vspace{0.3cm}
\textbf{Ý nghĩa:}
\begin{itemize}
  \item Hàm không được “nhảy đột ngột”; thay đổi trong đầu vào kéo theo thay đổi có kiểm soát ở đầu ra.
  \item Nếu đạo hàm của $h$ bị chặn bởi $M$ (tức $|\nabla h|\le M$), thì $h$ Lipschitz với hằng số $L\le M$.
\end{itemize}
\end{frame}
\begin{frame}{Áp dụng trong NTK}
\textbf{Áp dụng vào NTK:}
\[
g_{ir}(\theta_r)
= \nabla_{\theta_r} f_m(x_i;\theta)
\]
có dạng
\[
g_{ir} =
\frac{1}{\sqrt{m}}
\begin{pmatrix}
\sigma(w_r^\top x_i) \\
a_r \sigma'(w_r^\top x_i) x_i
\end{pmatrix}.
\]

Nhờ $\sigma$ trơn và $\sigma',\sigma''$ bị chặn:
\[
\|\nabla_{\theta_r} g_{ir}\| = O\!\left(\frac{1}{\sqrt{m}}\right)
\quad\Rightarrow\quad
g_{ir} \text{ là Lipschitz theo }\theta_r.
\]

Do đó:
\[
\|g_{ir}(t) - g_{ir}(0)\|
\le
O\!\left(\frac{1}{\sqrt{m}}\right)
\cdot
\|\theta_r(t)-\theta_r(0)\|
=
O\!\left(\frac{1}{m}\right).
\]
\end{frame}

\begin{frame}{Bước 2: thay đổi của một hạng tử NTK}

Đặt:
\[
  k_{ij}^{(r)}(t) := \langle g_{ir}(t),\, g_{jr}(t)\rangle.
\]

Hiệu giữa hai thời điểm:
\[
  |k_{ij}^{(r)}(t) - k_{ij}^{(r)}(0)|
  \le
  \|g_{ir}(t)-g_{ir}(0)\|\cdot \|g_{jr}(t)\|
  \;+\;
  \|g_{jr}(t)-g_{jr}(0)\|\cdot \|g_{ir}(0)\|.
\]

Do:
\[
  \|g_{ir}(t)-g_{ir}(0)\| = O(1/m), \qquad
  \|g_{jr}(t)\| = O(1/\sqrt{m}),
\]

ta được:
\[
  |k_{ij}^{(r)}(t) - k_{ij}^{(r)}(0)|
  = O\!\left(\frac{1}{m}\right)
    \cdot O\!\left(\frac{1}{\sqrt{m}}\right)
  = O\!\left(\frac{1}{m\sqrt{m}}\right).
\]

\end{frame}
\begin{frame}{Bước 2: tổng hợp tác động của toàn bộ $m$ neuron}

Tổng thay đổi của NTK:
\[
  |K_{m,ij}(t) - K_{m,ij}(0)|
  = \left|\sum_{r=1}^m \big(k_{ij}^{(r)}(t) - k_{ij}^{(r)}(0)\big)\right|.
\]

Mỗi hạng tử có bậc $O(1/(m\sqrt{m}))$, nên:
\[
  |K_{m,ij}(t) - K_{m,ij}(0)|
  \le
  m \cdot O\!\left(\frac{1}{m\sqrt{m}}\right)
  = O\!\left(\frac{1}{\sqrt{m}}\right)
  \xrightarrow[m\to\infty]{}
  0.
\]

\textbf{Kết luận Bước 2:}
\[
  \|K_m(\theta_m(t)) - K_m(\theta_m(0))\|
  = O\!\left(\frac{1}{\sqrt{m}}\right)
  \to 0.
\]

\end{frame}


%================= GOAL 2: SUMMARY =================
\begin{frame}{Mục tiêu 2 – Kết luận NTK regime}
  Gộp Bước 1 và Bước 2:
  \begin{itemize}
    \item (Bước 1) Tại khởi tạo,
      \[
        K_m(\theta_m(0)) \xrightarrow[m\to\infty]{a.s.} \Theta(X,X).
      \]
    \item (Bước 2) Trong quá trình train với thời gian hữu hạn,
      \[
        \sup_{t\in[0,T]}
        \|K_m(\theta_m(t)) - K_m(\theta_m(0))\|
        \xrightarrow[m\to\infty]{} 0.
      \]
  \end{itemize}

  \vspace{0.2cm}
  \textbf{Kết luận:} Trong NTK regime, ta có xấp xỉ:
  \[
    \frac{d\mathbf f(t)}{dt}
      = -K_m(\theta_m(t))(\mathbf f(t)-\mathbf y)
      \approx -\Theta(X,X)(\mathbf f(t)-\mathbf y).
  \]

\end{frame}

%================= GOAL 3 =================
\begin{frame}{Mục tiêu 3: Giải ODE và kernel regression}
  Từ Mục tiêu 2, ta xét hệ ODE xấp xỉ:
  \[
    \frac{d\mathbf f(t)}{dt}
    = -\Theta(X,X)\big(\mathbf f(t) - \mathbf y\big).
  \]
  Đặt $\mathbf u(t) = \mathbf f(t) - \mathbf y$:
  \[
    \frac{d\mathbf u(t)}{dt}
    = -\Theta(X,X)\mathbf u(t).
  \]
  Đây là ODE tuyến tính với nghiệm:
  \[
    \mathbf u(t)
    = e^{-\Theta(X,X)t}\mathbf u(0).
  \]
  Suy ra:
  \[
    \mathbf f(t)
    = \mathbf y + e^{-\Theta(X,X)t}(\mathbf f(0) - \mathbf y).
  \]

  Khi $t \to \infty$:
  \[
    e^{-\Theta(X,X)t} \to 0
    \Rightarrow
    \mathbf f_\infty = \mathbf y
    \quad \text{(fit đúng trên tập train)}.
  \]
\end{frame}

\begin{frame}{Giá trị tại điểm test $x$ và công thức NTK predictor}
  Với điểm test $x$, ta có động lực:
  \[
    \frac{df_t(x)}{dt}
    = -\Theta(x,X)\big(\mathbf f(t) - \mathbf y\big),
  \]
  với
  \(
    \Theta(x,X) = [\Theta(x,x_1),\dots,\Theta(x,x_n)].
  \)

  Thay biểu thức $\mathbf f(t)$ và giải, ta thu được nghiệm giới hạn:
  \[
    f_\infty(x)
    = f_0(x) - \Theta(x,X)\Theta(X,X)^{-1}(\mathbf f(0) - \mathbf y).
  \]

  Nếu khởi tạo sao cho $\mathbb E[f_0(x)]=0$ (hoặc bỏ qua term này), ta có:
  \[
    \boxed{
    f_\infty(x)
    = \Theta(x,X)\Theta(X,X)^{-1}\mathbf y
    }
  \]
  Đây chính là nghiệm của \textbf{kernel regression} với kernel $\Theta$.
\end{frame}

\begin{frame}{Tóm tắt toàn bộ chứng minh NTK}
  \begin{enumerate}
    \item \textbf{Mục tiêu 1:}  
          Từ gradient flow và Jacobian:
          \[
            \dot{\mathbf f}(t)
            = -K(\theta(t))(\mathbf f(t)-\mathbf y),
            \quad K = JJ^\top.
          \]
    \item \textbf{Mục tiêu 2:} (NTK regime)  
          Khi width $m\to\infty$:
          \[
            K_m(\theta(t)) \approx K_m(\theta(0)) \to \Theta(X,X),
          \]
          nên kernel gần như \textbf{hằng theo $t$}.
    \item \textbf{Mục tiêu 3:}  
          Thay $K$ bằng $\Theta(X,X)$, giải ODE:
          \[
            f_\infty(x)
            = \Theta(x,X)\Theta(X,X)^{-1}\mathbf y,
          \]
          $\Rightarrow$ gradient descent trên mạng rộng = kernel regression với NTK.
  \end{enumerate}
\end{frame}

\begin{frame}{Kết quả kinh điển: NNGP Kernel của ReLU}

Kết quả cổ điển từ Cho \& Saul (2009), Neal (1996), dùng lại trong
Jacot et al. (2018) và Lee et al. (2019):

\vspace{0.2cm}
Khi $w \sim \mathcal N(0,\tfrac{1}{k'}I)$, ta có kernel NNGP:
\[
K(x,x')
=
\frac{\|x\|\,\|x'\|}{2\pi k'}
\Big(
(\pi - \theta)\cos\theta + \sin\theta
\Big)
I_k
\]
trong đó
\[
\cos\theta
=
\frac{x^\top x'}{\|x\|\|x'\|}.
\]

\vspace{0.2cm}
Đây là \textbf{arc--cosine kernel cấp 1 (ReLU NNGP kernel)}.
\end{frame}

\begin{frame}{Kết quả kinh điển: NTK của mạng hai tầng ReLU}

Theo Jacot et al. (2018), Lee et al. (2019):
khi chiều rộng $n_h \to \infty$, NTK thực nghiệm hội tụ về
kernel xác định:

\[
\Theta(x,x')
=
\frac{x^\top x'}{2\pi k'}(\pi - \theta)
+
\frac{\|x\|\|x'\|}{2\pi k'}
\Big(
(\pi - \theta)\cos\theta + \sin\theta
\Big)
\; I_k
\]

với cùng định nghĩa
\[
\cos\theta = \frac{x^\top x'}{\|x\|\|x'\|}.
\]

\vspace{0.2cm}
NTK này gồm:
\begin{itemize}
    \item một phần từ đạo hàm theo $W_2$ (ReLU kernel),
    \item một phần từ đạo hàm theo $W_1$ (term $x^\top x'$).
\end{itemize}

\end{frame}

%================= PROBLEM & PIPELINE =================
\section{Vấn đề nghiên cứu & Pipeline}

\begin{frame}{Vấn đề cốt lõi: Algorithmic Generalization}
  \textbf{Câu hỏi nghiên cứu (Problem Statement):}
  \begin{itemize}
    \item Các mạng Neural truyền thống thường thất bại khi \textbf{ngoại suy (extrapolate)} độ dài. 
    \item Ví dụ: Train trên dãy số độ dài $K=5$, test trên độ dài $K=10$ thường sai hoàn toàn.
    \item \textbf{Mục tiêu:} Làm sao để một mạng Neural học được \textbf{chính xác thuật toán} (Exact Learning) để đảm bảo tính đúng đắn trên mọi độ dài?
  \end{itemize}

  \vspace{0.3cm}
  \textbf{Hướng tiếp cận của bài báo:}
  Sử dụng chế độ \textbf{Neural Tangent Kernel (NTK)} với mạng vô cùng rộng ($m \to \infty$) để đưa bài toán về dạng Kernel Regression tuyến tính.
\end{frame}

\begin{frame}{Pipeline của bài báo (The Pipeline)}
  Quy trình giải quyết vấn đề được chia thành các bước:
  \begin{enumerate}
    \item \textbf{Step 1: Định nghĩa Task:}
          Xác định hàm mục tiêu $f^*(x)$ là một thuật toán (ví dụ: hoán vị, cộng, nhân).
    \item \textbf{Step 2: Xây dựng Kernel $\Theta$:}
          Dùng kiến trúc mạng 2 tầng ReLU, đưa về NTK giới hạn $\Theta(X,X)$.
    \item \textbf{Step 3: Giải bài toán hồi quy (Optimization):}
          Nghiệm của mạng sau khi train vô hạn thời gian chính là nghiệm Kernel Ridge Regression:
          \[ f(x) = \Theta(x, X) \Theta(X, X)^{-1} Y \]
    \item \textbf{Step 4: Tìm biên dữ liệu (Sample Complexity Bound):}
          Xác định số lượng mẫu $N$ tối thiểu cần thiết để sai số $||f - f^*|| \approx 0$.
  \end{enumerate}
\end{frame}

%================= WHAT WAS PROVEN =================
\section{Điều Paper đã chứng minh (The Proof)}

\begin{frame}{Định lý Scaling Law (Điều đã chứng minh)}
  \textbf{Giả thuyết chính (The hypothesis):}
  Để học chính xác một thuật toán có độ phức tạp nhất định, kích thước tập huấn luyện $N$ phải tăng theo quy luật với độ dài đầu vào $K$.

  \vspace{0.3cm}
  \textbf{Kết quả chứng minh (Với bài toán Hoán vị):}
  Số lượng mẫu cần thiết $N$ tỷ lệ với:
  \[
    \boxed{ N \propto K^2 \log K }
  \]
  Nếu $N$ nhỏ hơn ngưỡng này, độ chính xác sẽ suy giảm đột ngột (Phase transition).
\end{frame}

\begin{frame}{Bằng chứng thực nghiệm (Validation)}
  \textbf{Mô tả thí nghiệm:}
  \begin{itemize}
    \item Chạy training trên Google Colab với mạng 50,000 neuron.
    \item Vary input size $K$ từ 2 đến 30.
    \item Đo $N$ tối thiểu để đạt Accuracy > 90\%.
  \end{itemize}

  \vspace{0.2cm}
  \textbf{Kết quả:} Đường thực nghiệm (Xanh dương) bám sát đường lý thuyết $O(K^2 \log K)$ (Xanh lá).
\end{frame}

\begin{frame}{Biểu đồ minh chứng (Evidence Plot)}
\centering
\resizebox{0.9\textwidth}{!}{
  \begin{tikzpicture}
    \begin{semilogyaxis}[
        title={Experiment 1: Ensemble size to achieve 90\% accuracy},
        xlabel={Input size $K$},
        ylabel={Ensemble size $N$},
        xmin=0, xmax=32,
        ymin=1, ymax=20000,
        xtick={2,4,6,8,10,12,14,16,18,20,22,24,26,28,30},
        legend pos=north west,
        ymajorgrids=true,
        grid style=dashed,
        width=12cm, height=8cm,
    ]
    % Theoretical Bound
    \addplot[color=green!60!black, style=solid, line width=1.5pt]
    coordinates {
        (2.0,5013.9) (3.5,5075.1) (4.9,5195.7) (6.4,5383.4) (7.9,5643.9) (9.4,5981.8) (10.8,6400.9) (12.3,6904.2) (13.8,7494.7) (15.3,8174.6) (16.7,8946.4) (18.2,9811.9) (19.7,10772.9) (21.2,11831.3) (22.6,12988.5) (24.1,14246.0) (25.6,15605.2) (27.1,17067.3) (28.5,18633.7) (30.0,20305.4) 
    };
    \addlegendentry{Lý thuyết $\mathcal{O}(K^2 \log K)$}
    
    % Empirical Data
    \addplot[color=blue!60!black, mark=*, style=solid, line width=1.5pt]
    coordinates {
        (2,1) (3,0) (4,39) (5,27) (6,16) (7,32) (8,39) (9,56) (10,58) (11,50) (12,51) (13,85) (14,86) (15,84) (16,84) (17,91) (18,117) (19,114) (20,131) (21,123) (22,123) (23,148) (24,137) (25,145) (26,167) (27,170) (28,182) (29,171) (30,171) 
    };
    \addlegendentry{Thực nghiệm (Colab 50k)}
    \end{semilogyaxis}
  \end{tikzpicture}
}
\end{frame}

%================= CONCLUSION =================
\section{Kết luận}
\begin{frame}{Kết luận chung}
  \begin{itemize}
    \item \textbf{Bài toán:} Đã giải quyết vấn đề "Học chính xác thuật toán" bằng NTK regime.
    \item \textbf{Pipeline:} Thiết lập pipeline rõ ràng từ Data $\to$ Kernel $\to$ Regression Solution.
    \item \textbf{Chứng minh:} Xác nhận thành công Scaling Law $N \propto K^2 \log K$.
    \item \textbf{Ý nghĩa:} Mở ra hướng đi mới cho việc áp dụng Neural Networks vào các bài toán suy luận thuật toán (Algorithmic Reasoning) đòi hỏi độ chính xác tuyệt đối.
  \end{itemize}
\end{frame}

\end{document}

